\documentclass[11pt]{article}
\usepackage[margin=1in]{geometry}
\usepackage{amsmath, amssymb, amsthm}
\usepackage{bm}

\title{\textbf{Continuous Relaxation for Graph Isomorphism Detection}}
\author{}
\date{}

\begin{document}
\maketitle

\section{Problem Formulation}

Given two simple, undirected graphs $G_1 = (V, E_1)$ and $G_2 = (V, E_2)$ on $n$ vertices, we seek to determine whether they are isomorphic. Two graphs are isomorphic if there exists a permutation $\pi: V \to V$ such that $(u, v) \in E_1 \iff (\pi(u), \pi(v)) \in E_2$.

Let $\bm{A}, \bm{B} \in \{0,1\}^{n \times n}$ denote the adjacency matrices of $G_1$ and $G_2$ respectively. A permutation $\pi$ can be represented as a permutation matrix $\bm{P} \in \{0,1\}^{n \times n}$ where $P_{ij} = 1$ iff $\pi(i) = j$. The graphs are isomorphic iff there exists a permutation matrix $\bm{P}$ such that:
\begin{equation}
    \bm{P}^\top \bm{A} \bm{P} = \bm{B}
\end{equation}

Since searching over all $n!$ permutations is computationally intractable, we relax the problem to a continuous optimization over \emph{doubly stochastic matrices}.

\section{Doubly Stochastic Relaxation}

A matrix $\bm{X} \in \mathbb{R}^{n \times n}$ is \textbf{doubly stochastic} if:
\begin{equation}
    X_{ij} \geq 0, \quad \sum_{j=1}^{n} X_{ij} = 1 \;\forall i, \quad \sum_{i=1}^{n} X_{ij} = 1 \;\forall j
\end{equation}

The set of doubly stochastic matrices forms the \emph{Birkhoff polytope} $\mathcal{B}_n$, whose vertices are exactly the $n!$ permutation matrices (Birkhoff--von Neumann theorem). By optimizing over $\mathcal{B}_n$, we obtain a convex relaxation of the combinatorial problem.

\section{Objective Function}

We formulate a quadratic program (QP) that penalizes structural mismatch between the graphs under a soft permutation $\bm{X} \in \mathcal{B}_n$. The implemented objective combines four signals:
\begin{enumerate}
    \item a degree-profile compatibility term,
    \item a first-order adjacency commutator penalty,
    \item a second-order (two-hop) commutator penalty,
    \item a spectral alignment penalty based on Laplacian eigenvectors.
\end{enumerate}

\subsection{Degree Profile Compatibility (Linear Term)}

Define the degree vectors
\begin{equation}
    d^A_i = \sum_k A_{ik},
    \qquad
    d^B_j = \sum_k B_{jk},
\end{equation}
and the one-hop degree-profile summaries
\begin{equation}
    s^A_i = \sum_k A_{ik} d^A_k,
    \qquad
    s^B_j = \sum_k B_{jk} d^B_k.
\end{equation}

The node-compatibility cost matrix is
\begin{equation}
    C_{ij} = \alpha \left| d^A_i - d^B_j \right|
    +
    \beta \left| s^A_i - s^B_j \right|,
\end{equation}
where $\alpha, \beta \ge 0$ are tunable weights. This penalizes mapping node $i$ in $G_1$ to node $j$ in $G_2$ when their local degree statistics differ. The linear objective term is:
\begin{equation}
    \mathcal{L}(\bm{X}) = \sum_{i=1}^{n} \sum_{j=1}^{n} C_{ij} X_{ij} = \langle \bm{C}, \bm{X} \rangle_F
\end{equation}

\subsection{Adjacency Mismatch (Quadratic Term)}

For a true permutation matrix $\bm{P}$, the transformed adjacency would be $\bm{P}^\top \bm{A} \bm{P}$. We penalize the discrepancy between $\bm{A}$ transformed by $\bm{X}$ and $\bm{B}$ transformed by $\bm{X}$:
\begin{equation}
    \mathcal{Q}(\bm{X}) = \sum_{p=1}^{n} \sum_{q=1}^{n} \left( \sum_{k=1}^{n} A_{pk} X_{kq} - \sum_{k=1}^{n} X_{pk} B_{kq} \right)^2
\end{equation}

This can be written in matrix form. Let $\bm{R} = \bm{A}\bm{X} - \bm{X}\bm{B}$, then:
\begin{equation}
    \mathcal{Q}(\bm{X}) = \|\bm{A}\bm{X} - \bm{X}\bm{B}\|_F^2 = \sum_{p,q} R_{pq}^2
\end{equation}

When $\bm{X} = \bm{P}$ is a valid isomorphism, we have $\bm{A}\bm{P} = \bm{P}\bm{B}$, hence $\mathcal{Q}(\bm{P}) = 0$.

\subsection{Two-Hop Neighborhood Penalty}

Direct adjacency alone does not fully capture local structure. Let
\begin{equation}
    \bm{A}^{(2)} = \bm{A}^2,
    \qquad
    \bm{B}^{(2)} = \bm{B}^2,
\end{equation}
whose entries count length-2 walks and therefore encode common-neighbor structure. We add the penalty
\begin{equation}
    \mathcal{T}(\bm{X}) = \|\bm{A}^{(2)}\bm{X} - \bm{X}\bm{B}^{(2)}\|_F^2.
\end{equation}

As with the first-order commutator, if $\bm{X}=\bm{P}$ is an exact isomorphism then $\bm{A}^2 \bm{P} = \bm{P}\bm{B}^2$, so $\mathcal{T}(\bm{P}) = 0$.

\subsection{Spectral Alignment Penalty}

Let the combinatorial Laplacians be
\begin{equation}
    \bm{L}_A = \operatorname{diag}(\bm{A}\bm{1}) - \bm{A},
    \qquad
    \bm{L}_B = \operatorname{diag}(\bm{B}\bm{1}) - \bm{B}.
\end{equation}
Let $\bm{U}_A, \bm{U}_B \in \mathbb{R}^{n \times n}$ be matrices of Laplacian eigenvectors ordered by increasing eigenvalue. In implementation, each eigenvector column is given a deterministic sign convention to reduce the usual sign ambiguity.

We then penalize disagreement between the spectral embedding of $G_1$ and the transported embedding of $G_2$:
\begin{equation}
    \mathcal{S}(\bm{X}) = \|\bm{U}_A - \bm{X}\bm{U}_B\|_F^2.
\end{equation}

This term gives the solver a global structural signal even when $\bm{X}$ is far from permutation-like.

\section{Quadratic Program}

The complete implemented optimization problem is:
\begin{equation}
\boxed{
\begin{aligned}
    Z^* = \min_{\bm{X}} \quad & \mathcal{L}(\bm{X})
    + \mathcal{Q}(\bm{X})
    + \lambda_T \mathcal{T}(\bm{X})
    + \lambda_S \mathcal{S}(\bm{X}) \\
    \text{s.t.} \quad & \sum_{j=1}^{n} X_{ij} = 1 \quad \forall\, i \in \{1, \ldots, n\} \\
    & \sum_{i=1}^{n} X_{ij} = 1 \quad \forall\, j \in \{1, \ldots, n\} \\
    & X_{ij} \geq 0 \quad \forall\, i, j
\end{aligned}
}
\end{equation}

Here $\lambda_T, \lambda_S \ge 0$ control the contribution of the two-hop and spectral penalties respectively. In the current codebase the default values are
\begin{equation}
    \alpha = 1.0, \qquad \beta = 0.35, \qquad \lambda_T = 0.2, \qquad \lambda_S = 0.15.
\end{equation}

The optimization remains a quadratic program over the Birkhoff polytope and is solved numerically with Gurobi.

\subsection{Remark on Entropy Regularization}

An additional entropy term,
\begin{equation}
    -\varepsilon \sum_{i,j} X_{ij}\log X_{ij},
\end{equation}
was proposed as a continuation-style regularizer to keep $\bm{X}$ diffuse early in the solve. This term is \emph{not} currently implemented in the codebase, because it requires a different nonlinear or piecewise-linear treatment than the present quadratic formulation.

\section{Isomorphism Index}

The optimal objective value $Z^*$ measures the ``distance'' from isomorphism:
\begin{itemize}
    \item $Z^* = 0$: The graphs are isomorphic (a perfect permutation exists)
    \item $Z^* > 0$: Structural mismatch; larger values indicate greater dissimilarity
\end{itemize}

We define a normalized \textbf{Isomorphism Index} $I \in (0, 1]$:
\begin{equation}
\boxed{
    I = \exp(-\lambda \cdot Z^*)
}
\end{equation}

where $\lambda > 0$ is a scaling parameter (default $\lambda = 0.1$). This maps the unbounded objective to a bounded similarity score:
\begin{itemize}
    \item $I \approx 1$: Graphs are (likely) isomorphic
    \item $I \ll 1$: Graphs are structurally dissimilar
\end{itemize}

\section{Interpretation of the Solution}

The optimal $\bm{X}^*$ is a doubly stochastic matrix representing a ``soft'' node correspondence:
\begin{itemize}
    \item If $\bm{X}^*$ is a permutation matrix (binary), it directly encodes the isomorphism mapping
    \item If $\bm{X}^*$ has fractional entries, the relaxation is not tight---this typically occurs for non-isomorphic graphs or for ambiguous graph symmetries
\end{itemize}

\section{Complexity}

The QP has $n^2$ variables and $2n$ equality constraints. Modern interior-point solvers achieve polynomial complexity in practice, making this approach tractable for moderate graph sizes ($n \lesssim 100$).

\end{document}
